%%%%%%%%%%%%%%%%%%%%%%%%%%%%%%%%%%%%%%%%%%%%%%%%%%%%%%%%%%%%
% 8.3 SAMPLING DISTRIBUTIONS
% The Composition of Advanced Mathematics
%%%%%%%%%%%%%%%%%%%%%%%%%%%%%%%%%%%%%%%%%%%%%%%%%%%%%%%%%%%%

\section{Sampling Distributions}
\label{sec:sampling-distributions}

\prerequisite{
Descriptive statistics, data collection and sampling, probability
distributions, random variables, expectation, variance, covariance,
binomial and normal distributions, laws of large numbers, central limit
theorems, and basic transformations of random variables.
}

\learningobjectives{
By the end of this section, the reader should be able to:
\begin{itemize}
    \item define a sampling distribution;
    \item distinguish a population distribution, sample distribution,
          and sampling distribution;
    \item explain why statistics are random variables before data are
          observed;
    \item derive the mean and variance of a sample mean;
    \item compute the standard error of a sample mean;
    \item apply the finite-population correction when appropriate;
    \item derive the sampling distribution of a sample proportion;
    \item compute the standard error of a sample proportion;
    \item use the central limit theorem for sample means;
    \item use normal approximations for sample proportions;
    \item identify exact normal sampling distributions;
    \item understand the distribution of sample sums;
    \item derive the chi-square distribution of normal sample variance;
    \item understand Student's \(t\)-distribution;
    \item understand the independence of sample mean and sample variance
          for normal data;
    \item understand the \(F\)-distribution as a ratio of scaled
          chi-square variables;
    \item derive sampling distributions for differences of means;
    \item derive sampling distributions for differences of proportions;
    \item understand paired-data sampling distributions;
    \item understand standard error versus standard deviation;
    \item interpret bias and variability of estimators;
    \item understand asymptotic normality;
    \item recognize the delta method conceptually;
    \item understand bootstrap sampling distributions conceptually;
    \item connect sampling distributions with estimation, confidence
          intervals, hypothesis testing, and regression.
\end{itemize}
}

%%%%%%%%%%%%%%%%%%%%%%%%%%%%%%%%%%%%%%%%%%%%%%%%%%%%%%%%%%%%
\subsection{What Is a Sampling Distribution?}
\label{subsec:what-is-sampling-distribution}
%%%%%%%%%%%%%%%%%%%%%%%%%%%%%%%%%%%%%%%%%%%%%%%%%%%%%%%%%%%%

Suppose a random sample is collected and a statistic

\[
T
=
T(X_1,\ldots,X_n)
\]

is computed.

Before the data are observed, the sample values

\[
X_1,\ldots,X_n
\]

are random variables.

Therefore \(T\) is also a random variable.

\begin{definitionbox}{Sampling Distribution}
The sampling distribution of a statistic is the probability distribution
of that statistic under repeated samples generated by the sampling
design or probability model.
\end{definitionbox}

Thus the statistic itself varies from sample to sample.

%%%%%%%%%%%%%%%%%%%%%%%%%%%%%%%%%%%%%%%%%%%%%%%%%%%%%%%%%%%%
\subsection{Population, Sample, and Sampling Distribution}
\label{subsec:population-sample-sampling-distribution}
%%%%%%%%%%%%%%%%%%%%%%%%%%%%%%%%%%%%%%%%%%%%%%%%%%%%%%%%%%%%

These three distributions must be distinguished carefully.

The \textbf{population distribution} describes individual population
values:

\[
\boxed{
X\sim F.
}
\]

The \textbf{sample distribution} describes the observations actually
seen in one realized sample:

\[
\boxed{
x_1,\ldots,x_n.
}
\]

The \textbf{sampling distribution} describes how a statistic would vary
across repeated samples:

\[
\boxed{
T(X_1,\ldots,X_n).
}
\]

%%%%%%%%%%%%%%%%%%%%%%%%%%%%%%%%%%%%%%%%%%%%%%%%%%%%%%%%%%%%
\subsection{Worked Example: Sampling Distribution from a Small Population}
\label{subsec:worked-small-population-sampling-distribution}
%%%%%%%%%%%%%%%%%%%%%%%%%%%%%%%%%%%%%%%%%%%%%%%%%%%%%%%%%%%%

\begin{workedexamplebox}{All Samples of Size Two}

Suppose the population is

\[
\{2,4,6\}.
\]

Take simple random samples of size \(2\) without replacement.

The possible samples are

\[
\{2,4\},
\qquad
\{2,6\},
\qquad
\{4,6\}.
\]

Their sample means are

\[
3,\qquad4,\qquad5.
\]

Each sample has probability

\[
\frac13.
\]

Therefore the sampling distribution of \(\overline{X}\) is

\[
\begin{array}{c|c}
\overline{x} & P(\overline{X}=\overline{x})\\
\hline
3 & 1/3\\
4 & 1/3\\
5 & 1/3
\end{array}
\]

\begin{answerbox}
\[
\boxed{
P(\overline{X}=3)
=
P(\overline{X}=4)
=
P(\overline{X}=5)
=
\frac13.
}
\]
\end{answerbox}

\end{workedexamplebox}

%%%%%%%%%%%%%%%%%%%%%%%%%%%%%%%%%%%%%%%%%%%%%%%%%%%%%%%%%%%%
\subsection{Statistics as Random Variables}
\label{subsec:statistics-random-variables}
%%%%%%%%%%%%%%%%%%%%%%%%%%%%%%%%%%%%%%%%%%%%%%%%%%%%%%%%%%%%

A statistic such as

\[
\overline{X}
=
\frac1n
\sum_{i=1}^{n}X_i
\]

is a function of random sample values.

Therefore \(\overline{X}\) has:

\[
\boxed{
\text{a mean},
}
\]

\[
\boxed{
\text{a variance},
}
\]

and

\[
\boxed{
\text{a probability distribution}.
}
\]

After the data are observed, the realized value

\[
\overline{x}
\]

is a number.

Before observation,

\[
\overline{X}
\]

is a random variable.

%%%%%%%%%%%%%%%%%%%%%%%%%%%%%%%%%%%%%%%%%%%%%%%%%%%%%%%%%%%%
\subsection{Estimator Versus Estimate}
\label{subsec:estimator-versus-estimate}
%%%%%%%%%%%%%%%%%%%%%%%%%%%%%%%%%%%%%%%%%%%%%%%%%%%%%%%%%%%%

An estimator is a random statistic used to estimate a parameter.

For example,

\[
\boxed{
\overline{X}
}
\]

is an estimator of

\[
\mu.
\]

Once the data are observed,

\[
\boxed{
\overline{x}
}
\]

is the numerical estimate.

Thus:

\[
\boxed{
\text{estimator}
=
\text{random quantity},
}
\]

while

\[
\boxed{
\text{estimate}
=
\text{observed numerical value}.
}
\]

%%%%%%%%%%%%%%%%%%%%%%%%%%%%%%%%%%%%%%%%%%%%%%%%%%%%%%%%%%%%
\subsection{Sampling Distribution of the Sample Mean}
\label{subsec:sampling-distribution-sample-mean}
%%%%%%%%%%%%%%%%%%%%%%%%%%%%%%%%%%%%%%%%%%%%%%%%%%%%%%%%%%%%

Let

\[
X_1,\ldots,X_n
\]

be independent and identically distributed random variables with

\[
\mathbb{E}[X_i]=\mu
\]

and

\[
\operatorname{Var}(X_i)=\sigma^2.
\]

The sample mean is

\[
\boxed{
\overline{X}
=
\frac1n
\sum_{i=1}^{n}X_i.
}
\]

Its expected value is

\[
\mathbb{E}[\overline{X}]
=
\frac1n
\sum_{i=1}^{n}
\mathbb{E}[X_i].
\]

Therefore,

\[
\boxed{
\mathbb{E}[\overline{X}]
=
\mu.
}
\]

%%%%%%%%%%%%%%%%%%%%%%%%%%%%%%%%%%%%%%%%%%%%%%%%%%%%%%%%%%%%
\subsection{Variance of the Sample Mean}
\label{subsec:variance-sample-mean}
%%%%%%%%%%%%%%%%%%%%%%%%%%%%%%%%%%%%%%%%%%%%%%%%%%%%%%%%%%%%

Using independence,

\[
\operatorname{Var}(\overline{X})
=
\operatorname{Var}
\left(
\frac1n
\sum_{i=1}^{n}X_i
\right).
\]

Hence,

\[
=
\frac1{n^2}
\sum_{i=1}^{n}
\operatorname{Var}(X_i).
\]

Thus,

\[
\boxed{
\operatorname{Var}(\overline{X})
=
\frac{\sigma^2}{n}.
}
\]

The standard deviation is

\[
\boxed{
\operatorname{SD}(\overline{X})
=
\frac{\sigma}{\sqrt{n}}.
}
\]

%%%%%%%%%%%%%%%%%%%%%%%%%%%%%%%%%%%%%%%%%%%%%%%%%%%%%%%%%%%%
\subsection{Standard Error}
\label{subsec:standard-error}
%%%%%%%%%%%%%%%%%%%%%%%%%%%%%%%%%%%%%%%%%%%%%%%%%%%%%%%%%%%%

\begin{definitionbox}{Standard Error}
The standard error of a statistic is the standard deviation of its
sampling distribution.
\end{definitionbox}

For the sample mean,

\[
\boxed{
\operatorname{SE}(\overline{X})
=
\frac{\sigma}{\sqrt{n}}.
}
\]

If \(\sigma\) is unknown, it is commonly estimated using \(s\):

\[
\boxed{
\widehat{\operatorname{SE}}(\overline{X})
=
\frac{s}{\sqrt{n}}.
}
\]

%%%%%%%%%%%%%%%%%%%%%%%%%%%%%%%%%%%%%%%%%%%%%%%%%%%%%%%%%%%%
\subsection{Standard Deviation Versus Standard Error}
\label{subsec:sd-versus-se}
%%%%%%%%%%%%%%%%%%%%%%%%%%%%%%%%%%%%%%%%%%%%%%%%%%%%%%%%%%%%

The population standard deviation

\[
\sigma
\]

describes variation among individual observations.

The standard error

\[
\frac{\sigma}{\sqrt{n}}
\]

describes variation among repeated sample means.

Thus:

\[
\boxed{
\text{SD}
=
\text{spread of observations},
}
\]

while

\[
\boxed{
\text{SE}
=
\text{spread of a statistic}.
}
\]

%%%%%%%%%%%%%%%%%%%%%%%%%%%%%%%%%%%%%%%%%%%%%%%%%%%%%%%%%%%%
\subsection{Worked Example: Standard Error of a Mean}
\label{subsec:worked-se-mean}
%%%%%%%%%%%%%%%%%%%%%%%%%%%%%%%%%%%%%%%%%%%%%%%%%%%%%%%%%%%%

\begin{workedexamplebox}{Sampling Mean Variability}

Suppose

\[
\mu=80,
\qquad
\sigma=12,
\qquad
n=36.
\]

Then

\[
\mathbb{E}[\overline{X}]
=
80.
\]

Also,

\[
\operatorname{SE}(\overline{X})
=
\frac{12}{6}.
\]

Therefore,

\begin{answerbox}
\[
\boxed{
\mathbb{E}[\overline{X}]=80,
\qquad
\operatorname{SE}(\overline{X})=2.
}
\]
\end{answerbox}

\end{workedexamplebox}

%%%%%%%%%%%%%%%%%%%%%%%%%%%%%%%%%%%%%%%%%%%%%%%%%%%%%%%%%%%%
\subsection{Effect of Sample Size}
\label{subsec:sample-size-se}
%%%%%%%%%%%%%%%%%%%%%%%%%%%%%%%%%%%%%%%%%%%%%%%%%%%%%%%%%%%%

Because

\[
\operatorname{SE}(\overline{X})
=
\frac{\sigma}{\sqrt{n}},
\]

larger samples produce more stable sample means.

If sample size is multiplied by \(4\),

\[
\boxed{
\operatorname{SE}
\text{ is divided by }2.
}
\]

If sample size is multiplied by \(9\),

\[
\boxed{
\operatorname{SE}
\text{ is divided by }3.
}
\]

This square-root relationship appears throughout statistics.

%%%%%%%%%%%%%%%%%%%%%%%%%%%%%%%%%%%%%%%%%%%%%%%%%%%%%%%%%%%%
\subsection{Finite-Population Sample Mean}
\label{subsec:finite-population-sample-mean}
%%%%%%%%%%%%%%%%%%%%%%%%%%%%%%%%%%%%%%%%%%%%%%%%%%%%%%%%%%%%

Suppose an SRS of size \(n\) is drawn without replacement from a finite
population of size \(N\).

Then the variance of the sample mean includes the finite-population
correction:

\[
\boxed{
\operatorname{Var}(\overline{X})
=
\frac{\sigma^2}{n}
\frac{N-n}{N-1}.
}
\]

Hence,

\[
\boxed{
\operatorname{SE}(\overline{X})
=
\frac{\sigma}{\sqrt{n}}
\sqrt{
\frac{N-n}{N-1}
}.
}
\]

%%%%%%%%%%%%%%%%%%%%%%%%%%%%%%%%%%%%%%%%%%%%%%%%%%%%%%%%%%%%
\subsection{Worked Example: Finite-Population Correction}
\label{subsec:worked-fpc-sampling-distribution}
%%%%%%%%%%%%%%%%%%%%%%%%%%%%%%%%%%%%%%%%%%%%%%%%%%%%%%%%%%%%

\begin{workedexamplebox}{Sampling a Large Fraction of a Population}

Suppose

\[
N=100,
\qquad
n=40,
\qquad
\sigma=15.
\]

The finite-population correction is

\[
\sqrt{
\frac{100-40}{99}
}
=
\sqrt{
\frac{60}{99}
}.
\]

Thus,

\[
\operatorname{SE}(\overline{X})
=
\frac{15}{\sqrt{40}}
\sqrt{
\frac{60}{99}
}.
\]

\begin{answerbox}
\[
\boxed{
\operatorname{SE}(\overline{X})
=
\frac{15}{\sqrt{40}}
\sqrt{
\frac{60}{99}
}.
}
\]
\end{answerbox}

\end{workedexamplebox}

%%%%%%%%%%%%%%%%%%%%%%%%%%%%%%%%%%%%%%%%%%%%%%%%%%%%%%%%%%%%
\subsection{Exact Normal Sampling Distribution}
\label{subsec:exact-normal-sampling-distribution}
%%%%%%%%%%%%%%%%%%%%%%%%%%%%%%%%%%%%%%%%%%%%%%%%%%%%%%%%%%%%

Suppose

\[
X_1,\ldots,X_n
\overset{\text{i.i.d.}}{\sim}
N(\mu,\sigma^2).
\]

Because linear combinations of independent normal variables are normal,

\[
\boxed{
\overline{X}
\sim
N
\left(
\mu,
\frac{\sigma^2}{n}
\right).
}
\]

This is exact for every positive integer \(n\).

No central limit approximation is needed.

%%%%%%%%%%%%%%%%%%%%%%%%%%%%%%%%%%%%%%%%%%%%%%%%%%%%%%%%%%%%
\subsection{Standardized Sample Mean with Known Variance}
\label{subsec:z-sample-mean-known-variance}
%%%%%%%%%%%%%%%%%%%%%%%%%%%%%%%%%%%%%%%%%%%%%%%%%%%%%%%%%%%%

If

\[
\overline{X}
\sim
N
\left(
\mu,
\frac{\sigma^2}{n}
\right),
\]

then

\[
\boxed{
Z
=
\frac{
\overline{X}-\mu
}{
\sigma/\sqrt{n}
}
\sim
N(0,1).
}
\]

This standardized form becomes the basis of \(z\)-based inference for
means when \(\sigma\) is known.

%%%%%%%%%%%%%%%%%%%%%%%%%%%%%%%%%%%%%%%%%%%%%%%%%%%%%%%%%%%%
\subsection{Central Limit Theorem for Sample Means}
\label{subsec:clt-sampling-distribution-mean}
%%%%%%%%%%%%%%%%%%%%%%%%%%%%%%%%%%%%%%%%%%%%%%%%%%%%%%%%%%%%

If the underlying population is not normal, the central limit theorem
still gives

\[
\boxed{
\frac{
\overline{X}-\mu
}{
\sigma/\sqrt{n}
}
\xrightarrow{d}
N(0,1)
}
\]

under standard i.i.d. finite-variance assumptions.

Therefore, for sufficiently large \(n\),

\[
\boxed{
\overline{X}
\approx
N
\left(
\mu,
\frac{\sigma^2}{n}
\right).
}
\]

%%%%%%%%%%%%%%%%%%%%%%%%%%%%%%%%%%%%%%%%%%%%%%%%%%%%%%%%%%%%
\subsection{Worked Example: CLT for a Sample Mean}
\label{subsec:worked-clt-sampling-mean}
%%%%%%%%%%%%%%%%%%%%%%%%%%%%%%%%%%%%%%%%%%%%%%%%%%%%%%%%%%%%

\begin{workedexamplebox}{Approximate Probability for a Sample Mean}

Suppose a population has

\[
\mu=100,
\qquad
\sigma=20.
\]

A random sample of size

\[
n=64
\]

is selected.

Then

\[
\operatorname{SE}(\overline{X})
=
\frac{20}{8}
=
2.5.
\]

Approximate

\[
P(\overline{X}>105).
\]

Standardizing,

\[
z
=
\frac{
105-100
}{
2.5
}
=
2.
\]

Therefore,

\begin{answerbox}
\[
\boxed{
P(\overline{X}>105)
\approx
1-\Phi(2).
}
\]
\end{answerbox}

\end{workedexamplebox}

%%%%%%%%%%%%%%%%%%%%%%%%%%%%%%%%%%%%%%%%%%%%%%%%%%%%%%%%%%%%
\subsection{Sampling Distribution of the Sample Sum}
\label{subsec:sampling-distribution-sample-sum}
%%%%%%%%%%%%%%%%%%%%%%%%%%%%%%%%%%%%%%%%%%%%%%%%%%%%%%%%%%%%

Define

\[
\boxed{
S_n
=
\sum_{i=1}^{n}X_i.
}
\]

For i.i.d. observations,

\[
\boxed{
\mathbb{E}[S_n]
=
n\mu,
}
\]

and

\[
\boxed{
\operatorname{Var}(S_n)
=
n\sigma^2.
}
\]

By the CLT,

\[
\boxed{
S_n
\approx
N(n\mu,n\sigma^2)
}
\]

for sufficiently large \(n\), under the usual assumptions.

%%%%%%%%%%%%%%%%%%%%%%%%%%%%%%%%%%%%%%%%%%%%%%%%%%%%%%%%%%%%
\subsection{Sample Proportion}
\label{subsec:sample-proportion-sampling-distribution}
%%%%%%%%%%%%%%%%%%%%%%%%%%%%%%%%%%%%%%%%%%%%%%%%%%%%%%%%%%%%

Suppose each observation records success or failure.

Let

\[
X_i
=
\begin{cases}
1,&\text{success},\\
0,&\text{failure}.
\end{cases}
\]

Then

\[
X_i\sim\operatorname{Bernoulli}(p).
\]

The sample proportion is

\[
\boxed{
\widehat{p}
=
\frac1n
\sum_{i=1}^{n}X_i.
}
\]

If

\[
Y
=
\sum_{i=1}^{n}X_i,
\]

then

\[
Y\sim\operatorname{Bin}(n,p)
\]

and

\[
\boxed{
\widehat{p}
=
\frac Yn.
}
\]

%%%%%%%%%%%%%%%%%%%%%%%%%%%%%%%%%%%%%%%%%%%%%%%%%%%%%%%%%%%%
\subsection{Mean of the Sample Proportion}
\label{subsec:mean-sample-proportion}
%%%%%%%%%%%%%%%%%%%%%%%%%%%%%%%%%%%%%%%%%%%%%%%%%%%%%%%%%%%%

Since

\[
\mathbb{E}[X_i]=p,
\]

we obtain

\[
\mathbb{E}[\widehat{p}]
=
\frac1n
\sum_{i=1}^{n}
p.
\]

Therefore,

\[
\boxed{
\mathbb{E}[\widehat{p}]
=
p.
}
\]

Thus the sample proportion is an unbiased estimator of the population
proportion under the Bernoulli sampling model.

%%%%%%%%%%%%%%%%%%%%%%%%%%%%%%%%%%%%%%%%%%%%%%%%%%%%%%%%%%%%
\subsection{Variance and Standard Error of a Proportion}
\label{subsec:variance-se-proportion}
%%%%%%%%%%%%%%%%%%%%%%%%%%%%%%%%%%%%%%%%%%%%%%%%%%%%%%%%%%%%

Because

\[
\operatorname{Var}(X_i)
=
p(1-p),
\]

we have

\[
\boxed{
\operatorname{Var}(\widehat{p})
=
\frac{
p(1-p)
}{n}.
}
\]

Thus,

\[
\boxed{
\operatorname{SE}(\widehat{p})
=
\sqrt{
\frac{
p(1-p)
}{n}
}.
}
\]

Since \(p\) is usually unknown, a common estimated standard error is

\[
\boxed{
\widehat{\operatorname{SE}}(\widehat{p})
=
\sqrt{
\frac{
\widehat{p}(1-\widehat{p})
}{n}
}.
}
\]

%%%%%%%%%%%%%%%%%%%%%%%%%%%%%%%%%%%%%%%%%%%%%%%%%%%%%%%%%%%%
\subsection{Normal Approximation for a Sample Proportion}
\label{subsec:normal-approx-sample-proportion}
%%%%%%%%%%%%%%%%%%%%%%%%%%%%%%%%%%%%%%%%%%%%%%%%%%%%%%%%%%%%

The central limit theorem gives

\[
\boxed{
\frac{
\widehat{p}-p
}{
\sqrt{
p(1-p)/n
}
}
\xrightarrow{d}
N(0,1).
}
\]

Hence,

\[
\boxed{
\widehat{p}
\approx
N
\left(
p,
\frac{
p(1-p)
}{n}
\right).
}
\]

A practical approximation generally requires both expected successes and
expected failures to be sufficiently large.

%%%%%%%%%%%%%%%%%%%%%%%%%%%%%%%%%%%%%%%%%%%%%%%%%%%%%%%%%%%%
\subsection{Worked Example: Sampling Distribution of a Proportion}
\label{subsec:worked-sampling-proportion}
%%%%%%%%%%%%%%%%%%%%%%%%%%%%%%%%%%%%%%%%%%%%%%%%%%%%%%%%%%%%

\begin{workedexamplebox}{Sampling Proportion}

Suppose

\[
p=0.40
\]

and

\[
n=200.
\]

Then

\[
\mathbb{E}[\widehat{p}]
=
0.40.
\]

The standard error is

\[
\sqrt{
\frac{
(0.40)(0.60)
}{200}
}.
\]

Therefore,

\begin{answerbox}
\[
\boxed{
\operatorname{SE}(\widehat{p})
=
\sqrt{
\frac{0.24}{200}
}.
}
\]
\end{answerbox}

\end{workedexamplebox}

%%%%%%%%%%%%%%%%%%%%%%%%%%%%%%%%%%%%%%%%%%%%%%%%%%%%%%%%%%%%
\subsection{Sample Proportion Without Replacement}
\label{subsec:sample-proportion-finite-population}
%%%%%%%%%%%%%%%%%%%%%%%%%%%%%%%%%%%%%%%%%%%%%%%%%%%%%%%%%%%%

Suppose a finite population contains \(N\) units, of which \(M\) are
successes.

Then

\[
p
=
\frac MN.
\]

If an SRS of size \(n\) is selected without replacement, the number of
successes has a hypergeometric distribution.

The sample proportion is

\[
\widehat{p}
=
\frac Yn,
\]

where \(Y\) is hypergeometric.

Its variance includes a finite-population correction.

%%%%%%%%%%%%%%%%%%%%%%%%%%%%%%%%%%%%%%%%%%%%%%%%%%%%%%%%%%%%
\subsection{Sampling Distribution of the Sample Variance}
\label{subsec:sampling-distribution-sample-variance}
%%%%%%%%%%%%%%%%%%%%%%%%%%%%%%%%%%%%%%%%%%%%%%%%%%%%%%%%%%%%

Suppose

\[
X_1,\ldots,X_n
\overset{\text{i.i.d.}}{\sim}
N(\mu,\sigma^2).
\]

The sample variance is

\[
S^2
=
\frac1{n-1}
\sum_{i=1}^{n}
(X_i-\overline{X})^2.
\]

A fundamental result is

\[
\boxed{
\frac{
(n-1)S^2
}{
\sigma^2
}
\sim
\chi^2_{n-1}.
}
\]

Thus normal sample variance is directly connected to the chi-square
distribution.

%%%%%%%%%%%%%%%%%%%%%%%%%%%%%%%%%%%%%%%%%%%%%%%%%%%%%%%%%%%%
\subsection{Chi-Square Distribution Review}
\label{subsec:chi-square-sampling-review}
%%%%%%%%%%%%%%%%%%%%%%%%%%%%%%%%%%%%%%%%%%%%%%%%%%%%%%%%%%%%

If

\[
Z_1,\ldots,Z_\nu
\]

are independent standard normal variables, then

\[
\boxed{
\sum_{i=1}^{\nu}Z_i^2
\sim
\chi^2_\nu.
}
\]

The parameter

\[
\nu
\]

is the number of degrees of freedom.

For the sample variance,

\[
\boxed{
\nu=n-1.
}
\]

%%%%%%%%%%%%%%%%%%%%%%%%%%%%%%%%%%%%%%%%%%%%%%%%%%%%%%%%%%%%
\subsection{Why the Degrees of Freedom Are \(n-1\)}
\label{subsec:variance-df-sampling}
%%%%%%%%%%%%%%%%%%%%%%%%%%%%%%%%%%%%%%%%%%%%%%%%%%%%%%%%%%%%

The deviations

\[
X_i-\overline{X}
\]

satisfy

\[
\sum_{i=1}^{n}
(X_i-\overline{X})
=
0.
\]

Therefore only

\[
n-1
\]

of them vary freely.

This loss of one degree of freedom appears explicitly in

\[
\chi^2_{n-1}.
\]

%%%%%%%%%%%%%%%%%%%%%%%%%%%%%%%%%%%%%%%%%%%%%%%%%%%%%%%%%%%%
\subsection{Mean and Variance of the Sample Variance}
\label{subsec:mean-variance-sample-variance}
%%%%%%%%%%%%%%%%%%%%%%%%%%%%%%%%%%%%%%%%%%%%%%%%%%%%%%%%%%%%

From

\[
\frac{
(n-1)S^2
}{
\sigma^2
}
\sim
\chi^2_{n-1},
\]

and

\[
\mathbb{E}[\chi^2_\nu]=\nu,
\]

we obtain

\[
\boxed{
\mathbb{E}[S^2]
=
\sigma^2.
}
\]

For normal data,

\[
\operatorname{Var}(\chi^2_\nu)=2\nu,
\]

so

\[
\boxed{
\operatorname{Var}(S^2)
=
\frac{
2\sigma^4
}{
n-1
}.
}
\]

%%%%%%%%%%%%%%%%%%%%%%%%%%%%%%%%%%%%%%%%%%%%%%%%%%%%%%%%%%%%
\subsection{Independence of Mean and Variance for Normal Samples}
\label{subsec:mean-variance-independence-normal}
%%%%%%%%%%%%%%%%%%%%%%%%%%%%%%%%%%%%%%%%%%%%%%%%%%%%%%%%%%%%

A remarkable property of normal samples is

\[
\boxed{
\overline{X}
\perp
S^2.
}
\]

That is, the sample mean and sample variance are independent.

This property does not generally hold for arbitrary population
distributions.

It plays a central role in deriving Student's \(t\)-distribution.

%%%%%%%%%%%%%%%%%%%%%%%%%%%%%%%%%%%%%%%%%%%%%%%%%%%%%%%%%%%%
\subsection{Student's \(t\)-Distribution}
\label{subsec:students-t-sampling}
%%%%%%%%%%%%%%%%%%%%%%%%%%%%%%%%%%%%%%%%%%%%%%%%%%%%%%%%%%%%

Suppose

\[
X_1,\ldots,X_n
\overset{\text{i.i.d.}}{\sim}
N(\mu,\sigma^2)
\]

but \(\sigma\) is unknown.

Then

\[
\boxed{
T
=
\frac{
\overline{X}-\mu
}{
S/\sqrt{n}
}
\sim
t_{n-1}.
}
\]

The distribution

\[
t_{n-1}
\]

is Student's \(t\)-distribution with \(n-1\) degrees of freedom.

%%%%%%%%%%%%%%%%%%%%%%%%%%%%%%%%%%%%%%%%%%%%%%%%%%%%%%%%%%%%
\subsection{Constructing the \(t\)-Distribution}
\label{subsec:t-construction-sampling}
%%%%%%%%%%%%%%%%%%%%%%%%%%%%%%%%%%%%%%%%%%%%%%%%%%%%%%%%%%%%

If

\[
Z\sim N(0,1)
\]

and

\[
U\sim\chi^2_\nu
\]

are independent, then

\[
\boxed{
T
=
\frac{
Z
}{
\sqrt{U/\nu}
}
\sim
t_\nu.
}
\]

For a normal sample,

\[
Z
=
\frac{
\overline{X}-\mu
}{
\sigma/\sqrt{n}
},
\]

and

\[
U
=
\frac{
(n-1)S^2
}{
\sigma^2
}.
\]

Substituting gives

\[
\boxed{
T
=
\frac{
\overline{X}-\mu
}{
S/\sqrt{n}
}.
}
\]

%%%%%%%%%%%%%%%%%%%%%%%%%%%%%%%%%%%%%%%%%%%%%%%%%%%%%%%%%%%%
\subsection{Shape of the \(t\)-Distribution}
\label{subsec:t-distribution-shape}
%%%%%%%%%%%%%%%%%%%%%%%%%%%%%%%%%%%%%%%%%%%%%%%%%%%%%%%%%%%%

The \(t\)-distribution is:

\[
\boxed{
\text{symmetric around zero},
}
\]

and has heavier tails than the standard normal.

The tail thickness depends on the degrees of freedom.

As

\[
\nu\to\infty,
\]

\[
\boxed{
t_\nu
\xrightarrow{d}
N(0,1).
}
\]

Thus the \(t\)-distribution approaches the standard normal as sample size
grows.

%%%%%%%%%%%%%%%%%%%%%%%%%%%%%%%%%%%%%%%%%%%%%%%%%%%%%%%%%%%%
\subsection{Why the \(t\)-Distribution Has Heavier Tails}
\label{subsec:why-t-heavy-tails}
%%%%%%%%%%%%%%%%%%%%%%%%%%%%%%%%%%%%%%%%%%%%%%%%%%%%%%%%%%%%

If \(\sigma\) were known, standardization would use

\[
\sigma.
\]

When \(\sigma\) is unknown, it is replaced by the random quantity

\[
S.
\]

This introduces additional uncertainty.

The heavier tails of the \(t\)-distribution account for that
uncertainty.

%%%%%%%%%%%%%%%%%%%%%%%%%%%%%%%%%%%%%%%%%%%%%%%%%%%%%%%%%%%%
\subsection{Worked Example: \(t\)-Statistic}
\label{subsec:worked-t-statistic}
%%%%%%%%%%%%%%%%%%%%%%%%%%%%%%%%%%%%%%%%%%%%%%%%%%%%%%%%%%%%

\begin{workedexamplebox}{Standardizing with an Estimated Standard Deviation}

Suppose

\[
n=25,
\qquad
\overline{x}=52,
\qquad
s=10,
\qquad
\mu_0=50.
\]

Then

\[
t
=
\frac{
52-50
}{
10/\sqrt{25}
}.
\]

Thus,

\[
t
=
\frac2{2}
=
1.
\]

\begin{answerbox}
\[
\boxed{
t=1,
\qquad
df=24.
}
\]
\end{answerbox}

\end{workedexamplebox}

%%%%%%%%%%%%%%%%%%%%%%%%%%%%%%%%%%%%%%%%%%%%%%%%%%%%%%%%%%%%
\subsection{\(F\)-Distribution}
\label{subsec:f-distribution-sampling}
%%%%%%%%%%%%%%%%%%%%%%%%%%%%%%%%%%%%%%%%%%%%%%%%%%%%%%%%%%%%

Suppose

\[
U_1\sim\chi^2_{\nu_1}
\]

and

\[
U_2\sim\chi^2_{\nu_2}
\]

are independent.

Then

\[
\boxed{
F
=
\frac{
U_1/\nu_1
}{
U_2/\nu_2
}
\sim
F_{\nu_1,\nu_2}.
}
\]

The \(F\)-distribution is supported on

\[
(0,\infty).
\]

It is central to comparisons of variances and analysis of variance.

%%%%%%%%%%%%%%%%%%%%%%%%%%%%%%%%%%%%%%%%%%%%%%%%%%%%%%%%%%%%
\subsection{Ratio of Two Normal Sample Variances}
\label{subsec:ratio-sample-variances}
%%%%%%%%%%%%%%%%%%%%%%%%%%%%%%%%%%%%%%%%%%%%%%%%%%%%%%%%%%%%

Suppose two independent normal samples satisfy

\[
S_1^2
\]

and

\[
S_2^2
\]

with population variances

\[
\sigma_1^2
\]

and

\[
\sigma_2^2.
\]

Then

\[
\boxed{
\frac{
S_1^2/\sigma_1^2
}{
S_2^2/\sigma_2^2
}
\sim
F_{n_1-1,n_2-1}.
}
\]

If

\[
\sigma_1^2=\sigma_2^2,
\]

then

\[
\boxed{
\frac{
S_1^2
}{
S_2^2
}
\sim
F_{n_1-1,n_2-1}.
}
\]

%%%%%%%%%%%%%%%%%%%%%%%%%%%%%%%%%%%%%%%%%%%%%%%%%%%%%%%%%%%%
\subsection{Sampling Distribution of a Difference of Means}
\label{subsec:difference-two-means-sampling}
%%%%%%%%%%%%%%%%%%%%%%%%%%%%%%%%%%%%%%%%%%%%%%%%%%%%%%%%%%%%

Suppose two independent samples satisfy

\[
X_1,\ldots,X_{n_1}
\]

and

\[
Y_1,\ldots,Y_{n_2},
\]

with

\[
\mathbb{E}[X_i]=\mu_X,
\qquad
\operatorname{Var}(X_i)=\sigma_X^2,
\]

and

\[
\mathbb{E}[Y_j]=\mu_Y,
\qquad
\operatorname{Var}(Y_j)=\sigma_Y^2.
\]

Then

\[
\boxed{
\mathbb{E}
[
\overline{X}-\overline{Y}
]
=
\mu_X-\mu_Y.
}
\]

%%%%%%%%%%%%%%%%%%%%%%%%%%%%%%%%%%%%%%%%%%%%%%%%%%%%%%%%%%%%
\subsection{Variance of a Difference of Independent Means}
\label{subsec:variance-difference-means}
%%%%%%%%%%%%%%%%%%%%%%%%%%%%%%%%%%%%%%%%%%%%%%%%%%%%%%%%%%%%

Because the samples are independent,

\[
\operatorname{Var}
(
\overline{X}-\overline{Y}
)
=
\operatorname{Var}(\overline{X})
+
\operatorname{Var}(\overline{Y}).
\]

Therefore,

\[
\boxed{
\operatorname{Var}
(
\overline{X}-\overline{Y}
)
=
\frac{
\sigma_X^2
}{
n_1
}
+
\frac{
\sigma_Y^2
}{
n_2
}.
}
\]

Hence,

\[
\boxed{
\operatorname{SE}
(
\overline{X}-\overline{Y}
)
=
\sqrt{
\frac{
\sigma_X^2
}{
n_1
}
+
\frac{
\sigma_Y^2
}{
n_2
}
}.
}
\]

%%%%%%%%%%%%%%%%%%%%%%%%%%%%%%%%%%%%%%%%%%%%%%%%%%%%%%%%%%%%
\subsection{Normal Sampling Distribution for Difference of Means}
\label{subsec:normal-difference-means}
%%%%%%%%%%%%%%%%%%%%%%%%%%%%%%%%%%%%%%%%%%%%%%%%%%%%%%%%%%%%

If both populations are normal, then

\[
\boxed{
\overline{X}-\overline{Y}
\sim
N
\left(
\mu_X-\mu_Y,
\frac{
\sigma_X^2
}{
n_1
}
+
\frac{
\sigma_Y^2
}{
n_2
}
\right).
}
\]

For nonnormal populations, the CLT provides an approximate version for
sufficiently large samples.

%%%%%%%%%%%%%%%%%%%%%%%%%%%%%%%%%%%%%%%%%%%%%%%%%%%%%%%%%%%%
\subsection{Worked Example: Difference of Means}
\label{subsec:worked-difference-means}
%%%%%%%%%%%%%%%%%%%%%%%%%%%%%%%%%%%%%%%%%%%%%%%%%%%%%%%%%%%%

\begin{workedexamplebox}{Comparing Two Independent Means}

Suppose

\[
\mu_X=70,
\qquad
\sigma_X=10,
\qquad
n_1=100,
\]

and

\[
\mu_Y=65,
\qquad
\sigma_Y=12,
\qquad
n_2=144.
\]

Then

\[
\mathbb{E}
[
\overline{X}-\overline{Y}
]
=
5.
\]

The standard error is

\[
\sqrt{
\frac{100}{100}
+
\frac{144}{144}
}.
\]

Therefore,

\begin{answerbox}
\[
\boxed{
\operatorname{SE}
(
\overline{X}-\overline{Y}
)
=
\sqrt{2}.
}
\]
\end{answerbox}

\end{workedexamplebox}

%%%%%%%%%%%%%%%%%%%%%%%%%%%%%%%%%%%%%%%%%%%%%%%%%%%%%%%%%%%%
\subsection{Pooled Variance Preview}
\label{subsec:pooled-variance-preview}
%%%%%%%%%%%%%%%%%%%%%%%%%%%%%%%%%%%%%%%%%%%%%%%%%%%%%%%%%%%%

If two independent normal populations have the same unknown variance,

\[
\sigma_1^2=\sigma_2^2=\sigma^2,
\]

the two sample variances may be combined.

The pooled estimator is

\[
\boxed{
S_p^2
=
\frac{
(n_1-1)S_1^2
+
(n_2-1)S_2^2
}{
n_1+n_2-2
}.
}
\]

This leads to the classical pooled two-sample \(t\)-distribution.

%%%%%%%%%%%%%%%%%%%%%%%%%%%%%%%%%%%%%%%%%%%%%%%%%%%%%%%%%%%%
\subsection{Welch Sampling Distribution Preview}
\label{subsec:welch-sampling-preview}
%%%%%%%%%%%%%%%%%%%%%%%%%%%%%%%%%%%%%%%%%%%%%%%%%%%%%%%%%%%%

When the population variances need not be equal, a standardized
difference of sample means uses

\[
\boxed{
\frac{
\overline{X}-\overline{Y}
-
(\mu_X-\mu_Y)
}{
\sqrt{
S_X^2/n_1
+
S_Y^2/n_2
}
}.
}
\]

Its exact distribution is not generally a simple \(t\)-distribution.

Welch's method approximates it by a \(t\)-distribution with estimated
degrees of freedom.

This is developed later in hypothesis testing and interval estimation.

%%%%%%%%%%%%%%%%%%%%%%%%%%%%%%%%%%%%%%%%%%%%%%%%%%%%%%%%%%%%
\subsection{Sampling Distribution of a Difference of Proportions}
\label{subsec:difference-proportions-sampling}
%%%%%%%%%%%%%%%%%%%%%%%%%%%%%%%%%%%%%%%%%%%%%%%%%%%%%%%%%%%%

Suppose two independent samples produce proportions

\[
\widehat{p}_1
\]

and

\[
\widehat{p}_2.
\]

Then

\[
\boxed{
\mathbb{E}
[
\widehat{p}_1-\widehat{p}_2
]
=
p_1-p_2.
}
\]

Because the samples are independent,

\[
\boxed{
\operatorname{Var}
(
\widehat{p}_1-\widehat{p}_2
)
=
\frac{
p_1(1-p_1)
}{
n_1
}
+
\frac{
p_2(1-p_2)
}{
n_2
}.
}
\]

%%%%%%%%%%%%%%%%%%%%%%%%%%%%%%%%%%%%%%%%%%%%%%%%%%%%%%%%%%%%
\subsection{Standard Error of Difference of Proportions}
\label{subsec:se-difference-proportions}
%%%%%%%%%%%%%%%%%%%%%%%%%%%%%%%%%%%%%%%%%%%%%%%%%%%%%%%%%%%%

The standard error is

\[
\boxed{
\operatorname{SE}
(
\widehat{p}_1-\widehat{p}_2
)
=
\sqrt{
\frac{
p_1(1-p_1)
}{
n_1
}
+
\frac{
p_2(1-p_2)
}{
n_2
}
}.
}
\]

In applications, the unknown population proportions are usually replaced
by sample-based estimates, depending on the inferential procedure.

%%%%%%%%%%%%%%%%%%%%%%%%%%%%%%%%%%%%%%%%%%%%%%%%%%%%%%%%%%%%
\subsection{Normal Approximation for Difference of Proportions}
\label{subsec:normal-difference-proportions}
%%%%%%%%%%%%%%%%%%%%%%%%%%%%%%%%%%%%%%%%%%%%%%%%%%%%%%%%%%%%

For sufficiently large independent samples,

\[
\boxed{
\widehat{p}_1-\widehat{p}_2
\approx
N
\left(
p_1-p_2,
\frac{
p_1(1-p_1)
}{
n_1
}
+
\frac{
p_2(1-p_2)
}{
n_2
}
\right).
}
\]

This becomes the basis of two-proportion confidence intervals and tests.

%%%%%%%%%%%%%%%%%%%%%%%%%%%%%%%%%%%%%%%%%%%%%%%%%%%%%%%%%%%%
\subsection{Worked Example: Difference of Proportions}
\label{subsec:worked-difference-proportions}
%%%%%%%%%%%%%%%%%%%%%%%%%%%%%%%%%%%%%%%%%%%%%%%%%%%%%%%%%%%%

\begin{workedexamplebox}{Two Population Proportions}

Suppose

\[
p_1=0.60,
\qquad
n_1=200,
\]

and

\[
p_2=0.50,
\qquad
n_2=300.
\]

Then

\[
\mathbb{E}
[
\widehat{p}_1-\widehat{p}_2
]
=
0.10.
\]

The standard error is

\[
\sqrt{
\frac{
(0.60)(0.40)
}{200}
+
\frac{
(0.50)(0.50)
}{300}
}.
\]

\begin{answerbox}
\[
\boxed{
\operatorname{SE}
(
\widehat{p}_1-\widehat{p}_2
)
=
\sqrt{
\frac{0.24}{200}
+
\frac{0.25}{300}
}.
}
\]
\end{answerbox}

\end{workedexamplebox}

%%%%%%%%%%%%%%%%%%%%%%%%%%%%%%%%%%%%%%%%%%%%%%%%%%%%%%%%%%%%
\subsection{Paired Sampling}
\label{subsec:paired-sampling-distribution}
%%%%%%%%%%%%%%%%%%%%%%%%%%%%%%%%%%%%%%%%%%%%%%%%%%%%%%%%%%%%

Suppose observations naturally occur in pairs:

\[
(X_1,Y_1),
\ldots,
(X_n,Y_n).
\]

Define differences

\[
\boxed{
D_i
=
X_i-Y_i.
}
\]

Then the paired problem becomes a one-sample problem involving

\[
D_1,\ldots,D_n.
\]

The sample mean difference is

\[
\boxed{
\overline{D}
=
\frac1n
\sum_{i=1}^{n}D_i.
}
\]

%%%%%%%%%%%%%%%%%%%%%%%%%%%%%%%%%%%%%%%%%%%%%%%%%%%%%%%%%%%%
\subsection{Sampling Distribution of Paired Mean Difference}
\label{subsec:paired-mean-sampling}
%%%%%%%%%%%%%%%%%%%%%%%%%%%%%%%%%%%%%%%%%%%%%%%%%%%%%%%%%%%%

If the differences are i.i.d. with

\[
\mathbb{E}[D_i]=\mu_D
\]

and

\[
\operatorname{Var}(D_i)=\sigma_D^2,
\]

then

\[
\boxed{
\mathbb{E}[\overline{D}]
=
\mu_D,
}
\]

and

\[
\boxed{
\operatorname{SE}(\overline{D})
=
\frac{
\sigma_D
}{
\sqrt{n}
}.
}
\]

For normal differences,

\[
\boxed{
\frac{
\overline{D}-\mu_D
}{
S_D/\sqrt{n}
}
\sim
t_{n-1}.
}
\]

%%%%%%%%%%%%%%%%%%%%%%%%%%%%%%%%%%%%%%%%%%%%%%%%%%%%%%%%%%%%
\subsection{Why Paired and Independent Samples Differ}
\label{subsec:paired-versus-independent}
%%%%%%%%%%%%%%%%%%%%%%%%%%%%%%%%%%%%%%%%%%%%%%%%%%%%%%%%%%%%

For paired observations,

\[
X_i
\]

and

\[
Y_i
\]

within the same pair may be strongly dependent.

The relevant variance is

\[
\operatorname{Var}(X_i-Y_i).
\]

Using covariance,

\[
\boxed{
\operatorname{Var}(X-Y)
=
\sigma_X^2
+
\sigma_Y^2
-
2\operatorname{Cov}(X,Y).
}
\]

Positive within-pair correlation can substantially reduce the variance of
the difference.

%%%%%%%%%%%%%%%%%%%%%%%%%%%%%%%%%%%%%%%%%%%%%%%%%%%%%%%%%%%%
\subsection{Worked Example: Benefit of Pairing}
\label{subsec:worked-benefit-pairing}
%%%%%%%%%%%%%%%%%%%%%%%%%%%%%%%%%%%%%%%%%%%%%%%%%%%%%%%%%%%%

\begin{workedexamplebox}{Variance of a Paired Difference}

Suppose

\[
\sigma_X=10,
\qquad
\sigma_Y=10,
\]

and

\[
\operatorname{Corr}(X,Y)=0.8.
\]

Then

\[
\operatorname{Cov}(X,Y)
=
(0.8)(10)(10)
=
80.
\]

Therefore,

\[
\operatorname{Var}(X-Y)
=
100+100-160.
\]

Thus,

\begin{answerbox}
\[
\boxed{
\operatorname{Var}(X-Y)=40.
}
\]
\end{answerbox}

Without the positive pairing covariance, the variance would have been
\(200\).

\end{workedexamplebox}

%%%%%%%%%%%%%%%%%%%%%%%%%%%%%%%%%%%%%%%%%%%%%%%%%%%%%%%%%%%%
\subsection{Bias of an Estimator}
\label{subsec:bias-estimator}
%%%%%%%%%%%%%%%%%%%%%%%%%%%%%%%%%%%%%%%%%%%%%%%%%%%%%%%%%%%%

\begin{definitionbox}{Bias}
The bias of an estimator \(\widehat{\theta}\) for parameter \(\theta\) is

\[
\boxed{
\operatorname{Bias}(\widehat{\theta})
=
\mathbb{E}[\widehat{\theta}]
-
\theta.
}
\]
\end{definitionbox}

If

\[
\mathbb{E}[\widehat{\theta}]
=
\theta,
\]

the estimator is unbiased.

%%%%%%%%%%%%%%%%%%%%%%%%%%%%%%%%%%%%%%%%%%%%%%%%%%%%%%%%%%%%
\subsection{Unbiased Estimators}
\label{subsec:unbiased-estimators}
%%%%%%%%%%%%%%%%%%%%%%%%%%%%%%%%%%%%%%%%%%%%%%%%%%%%%%%%%%%%

For i.i.d. sampling,

\[
\boxed{
\mathbb{E}[\overline{X}]
=
\mu,
}
\]

so the sample mean is unbiased for \(\mu\).

Likewise,

\[
\boxed{
\mathbb{E}[S^2]
=
\sigma^2,
}
\]

so the usual sample variance is unbiased for \(\sigma^2\).

However,

\[
S
\]

itself is generally not exactly unbiased for \(\sigma\).

%%%%%%%%%%%%%%%%%%%%%%%%%%%%%%%%%%%%%%%%%%%%%%%%%%%%%%%%%%%%
\subsection{Variance of an Estimator}
\label{subsec:variance-estimator}
%%%%%%%%%%%%%%%%%%%%%%%%%%%%%%%%%%%%%%%%%%%%%%%%%%%%%%%%%%%%

Two unbiased estimators may have different variances.

The estimator with smaller variance is more concentrated around the true
parameter.

Thus good estimators are often judged by both:

\[
\boxed{
\text{bias}
}
\]

and

\[
\boxed{
\text{variance}.
}
\]

%%%%%%%%%%%%%%%%%%%%%%%%%%%%%%%%%%%%%%%%%%%%%%%%%%%%%%%%%%%%
\subsection{Mean Squared Error}
\label{subsec:mean-squared-error}
%%%%%%%%%%%%%%%%%%%%%%%%%%%%%%%%%%%%%%%%%%%%%%%%%%%%%%%%%%%%

\begin{definitionbox}{Mean Squared Error}
The mean squared error of an estimator \(\widehat{\theta}\) is

\[
\boxed{
\operatorname{MSE}(\widehat{\theta})
=
\mathbb{E}
\left[
(\widehat{\theta}-\theta)^2
\right].
}
\]
\end{definitionbox}

It decomposes as

\[
\boxed{
\operatorname{MSE}(\widehat{\theta})
=
\operatorname{Var}(\widehat{\theta})
+
\operatorname{Bias}(\widehat{\theta})^2.
}
\]

%%%%%%%%%%%%%%%%%%%%%%%%%%%%%%%%%%%%%%%%%%%%%%%%%%%%%%%%%%%%
\subsection{Deriving the MSE Decomposition}
\label{subsec:mse-decomposition}
%%%%%%%%%%%%%%%%%%%%%%%%%%%%%%%%%%%%%%%%%%%%%%%%%%%%%%%%%%%%

Write

\[
\widehat{\theta}-\theta
=
\left(
\widehat{\theta}
-
\mathbb{E}[\widehat{\theta}]
\right)
+
\left(
\mathbb{E}[\widehat{\theta}]
-\theta
\right).
\]

Squaring and taking expectations gives

\[
\boxed{
\operatorname{MSE}
=
\operatorname{Variance}
+
\operatorname{Bias}^2.
}
\]

The cross term vanishes because

\[
\mathbb{E}
[
\widehat{\theta}
-
\mathbb{E}[\widehat{\theta}]
]
=
0.
\]

%%%%%%%%%%%%%%%%%%%%%%%%%%%%%%%%%%%%%%%%%%%%%%%%%%%%%%%%%%%%
\subsection{Consistency}
\label{subsec:consistency-sampling}
%%%%%%%%%%%%%%%%%%%%%%%%%%%%%%%%%%%%%%%%%%%%%%%%%%%%%%%%%%%%

An estimator sequence

\[
\widehat{\theta}_n
\]

is consistent for \(\theta\) if

\[
\boxed{
\widehat{\theta}_n
\xrightarrow{P}
\theta.
}
\]

For example, under the law of large numbers,

\[
\boxed{
\overline{X}_n
\xrightarrow{P}
\mu.
}
\]

Thus the sample mean is consistent for the population mean under standard
conditions.

%%%%%%%%%%%%%%%%%%%%%%%%%%%%%%%%%%%%%%%%%%%%%%%%%%%%%%%%%%%%
\subsection{Unbiasedness Versus Consistency}
\label{subsec:unbiasedness-versus-consistency}
%%%%%%%%%%%%%%%%%%%%%%%%%%%%%%%%%%%%%%%%%%%%%%%%%%%%%%%%%%%%

Unbiasedness concerns the estimator's expectation at a fixed sample size:

\[
\mathbb{E}[\widehat{\theta}_n]=\theta.
\]

Consistency concerns behavior as

\[
n\to\infty.
\]

An estimator can be biased but consistent.

An unbiased estimator need not automatically have small variance.

Thus these are different properties.

%%%%%%%%%%%%%%%%%%%%%%%%%%%%%%%%%%%%%%%%%%%%%%%%%%%%%%%%%%%%
\subsection{Asymptotic Distribution}
\label{subsec:asymptotic-distribution}
%%%%%%%%%%%%%%%%%%%%%%%%%%%%%%%%%%%%%%%%%%%%%%%%%%%%%%%%%%%%

An estimator may not have a simple exact finite-sample distribution.

Instead, after suitable scaling,

\[
a_n
(
\widehat{\theta}_n-\theta
)
\]

may converge in distribution.

A common result is

\[
\boxed{
\sqrt{n}
(
\widehat{\theta}_n-\theta
)
\xrightarrow{d}
N(0,V).
}
\]

This is called asymptotic normality.

%%%%%%%%%%%%%%%%%%%%%%%%%%%%%%%%%%%%%%%%%%%%%%%%%%%%%%%%%%%%
\subsection{Asymptotic Standard Error}
\label{subsec:asymptotic-standard-error}
%%%%%%%%%%%%%%%%%%%%%%%%%%%%%%%%%%%%%%%%%%%%%%%%%%%%%%%%%%%%

If

\[
\sqrt{n}
(
\widehat{\theta}_n-\theta
)
\xrightarrow{d}
N(0,V),
\]

then approximately

\[
\widehat{\theta}_n
\sim
N
\left(
\theta,
\frac{V}{n}
\right).
\]

Thus the asymptotic standard error is approximately

\[
\boxed{
\sqrt{
\frac{V}{n}
}.
}
\]

The quantity \(V\) may depend on unknown population parameters and must
often be estimated.

%%%%%%%%%%%%%%%%%%%%%%%%%%%%%%%%%%%%%%%%%%%%%%%%%%%%%%%%%%%%
\subsection{Delta Method Preview}
\label{subsec:delta-method-sampling}
%%%%%%%%%%%%%%%%%%%%%%%%%%%%%%%%%%%%%%%%%%%%%%%%%%%%%%%%%%%%

Suppose

\[
\sqrt{n}
(
\widehat{\theta}_n-\theta
)
\xrightarrow{d}
N(0,V)
\]

and \(g\) is differentiable at \(\theta\).

Then

\[
\boxed{
\sqrt{n}
\left[
g(\widehat{\theta}_n)
-
g(\theta)
\right]
\xrightarrow{d}
N
\left(
0,
[g'(\theta)]^2V
\right).
}
\]

This is the delta method.

It transfers asymptotic normality through smooth transformations.

%%%%%%%%%%%%%%%%%%%%%%%%%%%%%%%%%%%%%%%%%%%%%%%%%%%%%%%%%%%%
\subsection{Worked Example: Delta Method for a Logarithm}
\label{subsec:worked-delta-log}
%%%%%%%%%%%%%%%%%%%%%%%%%%%%%%%%%%%%%%%%%%%%%%%%%%%%%%%%%%%%

\begin{workedexamplebox}{Transforming an Asymptotically Normal Estimator}

Suppose

\[
\sqrt{n}
(
\widehat{\theta}-\theta
)
\xrightarrow{d}
N(0,V),
\qquad
\theta>0.
\]

Let

\[
g(\theta)=\ln\theta.
\]

Then

\[
g'(\theta)=\frac1\theta.
\]

Therefore,

\begin{answerbox}
\[
\boxed{
\sqrt{n}
\left(
\ln\widehat{\theta}
-
\ln\theta
\right)
\xrightarrow{d}
N
\left(
0,
\frac{V}{\theta^2}
\right).
}
\]
\end{answerbox}

\end{workedexamplebox}

%%%%%%%%%%%%%%%%%%%%%%%%%%%%%%%%%%%%%%%%%%%%%%%%%%%%%%%%%%%%
\subsection{Sampling Distribution of a Linear Combination}
\label{subsec:sampling-linear-combination}
%%%%%%%%%%%%%%%%%%%%%%%%%%%%%%%%%%%%%%%%%%%%%%%%%%%%%%%%%%%%

Suppose

\[
T
=
a_1X_1+\cdots+a_nX_n.
\]

Then

\[
\boxed{
\mathbb{E}[T]
=
\sum_{i=1}^{n}
a_i\mathbb{E}[X_i].
}
\]

If the variables are independent,

\[
\boxed{
\operatorname{Var}(T)
=
\sum_{i=1}^{n}
a_i^2
\operatorname{Var}(X_i).
}
\]

This framework includes sample means, weighted means, and contrasts.

%%%%%%%%%%%%%%%%%%%%%%%%%%%%%%%%%%%%%%%%%%%%%%%%%%%%%%%%%%%%
\subsection{Sampling Distribution of a Weighted Mean}
\label{subsec:sampling-weighted-mean}
%%%%%%%%%%%%%%%%%%%%%%%%%%%%%%%%%%%%%%%%%%%%%%%%%%%%%%%%%%%%

Suppose

\[
\widehat{\mu}
=
\sum_{i=1}^{n}
w_iX_i,
\]

where

\[
\sum_{i=1}^{n}w_i=1.
\]

If each

\[
\mathbb{E}[X_i]=\mu,
\]

then

\[
\boxed{
\mathbb{E}[\widehat{\mu}]
=
\mu.
}
\]

If the observations are independent,

\[
\boxed{
\operatorname{Var}(\widehat{\mu})
=
\sum_{i=1}^{n}
w_i^2\sigma_i^2.
}
\]

%%%%%%%%%%%%%%%%%%%%%%%%%%%%%%%%%%%%%%%%%%%%%%%%%%%%%%%%%%%%
\subsection{Order Statistics Sampling Distributions}
\label{subsec:order-statistics-sampling}
%%%%%%%%%%%%%%%%%%%%%%%%%%%%%%%%%%%%%%%%%%%%%%%%%%%%%%%%%%%%

For an i.i.d. sample, let

\[
X_{(1)}
\leq
\cdots
\leq
X_{(n)}
\]

denote the order statistics.

Their distributions are sampling distributions because each
\(X_{(k)}\) is a statistic.

If the population has CDF \(F\) and density \(f\), then

\[
\boxed{
f_{X_{(k)}}(x)
=
\frac{
n!
}{
(k-1)!(n-k)!
}
[F(x)]^{k-1}
[1-F(x)]^{n-k}
f(x).
}
\]

%%%%%%%%%%%%%%%%%%%%%%%%%%%%%%%%%%%%%%%%%%%%%%%%%%%%%%%%%%%%
\subsection{Sampling Distribution of the Minimum}
\label{subsec:sample-minimum-distribution}
%%%%%%%%%%%%%%%%%%%%%%%%%%%%%%%%%%%%%%%%%%%%%%%%%%%%%%%%%%%%

For

\[
X_{(1)}
=
\min(X_1,\ldots,X_n),
\]

\[
P(X_{(1)}>x)
=
P(X_1>x,\ldots,X_n>x).
\]

Under independence,

\[
=
[1-F(x)]^n.
\]

Therefore,

\[
\boxed{
F_{X_{(1)}}(x)
=
1-[1-F(x)]^n.
}
\]

%%%%%%%%%%%%%%%%%%%%%%%%%%%%%%%%%%%%%%%%%%%%%%%%%%%%%%%%%%%%
\subsection{Sampling Distribution of the Maximum}
\label{subsec:sample-maximum-distribution}
%%%%%%%%%%%%%%%%%%%%%%%%%%%%%%%%%%%%%%%%%%%%%%%%%%%%%%%%%%%%

For

\[
X_{(n)}
=
\max(X_1,\ldots,X_n),
\]

\[
P(X_{(n)}\leq x)
=
P(X_1\leq x,\ldots,X_n\leq x).
\]

Hence,

\[
\boxed{
F_{X_{(n)}}(x)
=
[F(x)]^n.
}
\]

%%%%%%%%%%%%%%%%%%%%%%%%%%%%%%%%%%%%%%%%%%%%%%%%%%%%%%%%%%%%
\subsection{Sampling Distribution of the Median Preview}
\label{subsec:sample-median-distribution-preview}
%%%%%%%%%%%%%%%%%%%%%%%%%%%%%%%%%%%%%%%%%%%%%%%%%%%%%%%%%%%%

The sample median is an order statistic or an average of two central
order statistics, depending on sample size.

Its exact sampling distribution can be derived from order-statistic
theory.

For large samples, under regularity conditions,

\[
\boxed{
\sqrt{n}
(
\widetilde{X}-m
)
\xrightarrow{d}
N
\left(
0,
\frac{
1
}{
4f(m)^2
}
\right),
}
\]

where \(m\) is the population median and \(f(m)>0\).

%%%%%%%%%%%%%%%%%%%%%%%%%%%%%%%%%%%%%%%%%%%%%%%%%%%%%%%%%%%%
\subsection{Bootstrap Sampling Distribution}
\label{subsec:bootstrap-sampling-distribution}
%%%%%%%%%%%%%%%%%%%%%%%%%%%%%%%%%%%%%%%%%%%%%%%%%%%%%%%%%%%%

When an analytic sampling distribution is difficult to derive, it may be
approximated computationally.

The bootstrap treats the observed empirical distribution as an estimate
of the population distribution.

A bootstrap sample is generated by sampling

\[
n
\]

observations with replacement from the observed data.

%%%%%%%%%%%%%%%%%%%%%%%%%%%%%%%%%%%%%%%%%%%%%%%%%%%%%%%%%%%%
\subsection{Bootstrap Algorithm}
\label{subsec:bootstrap-algorithm}
%%%%%%%%%%%%%%%%%%%%%%%%%%%%%%%%%%%%%%%%%%%%%%%%%%%%%%%%%%%%

Given observed data

\[
x_1,\ldots,x_n,
\]

a basic nonparametric bootstrap proceeds as follows:

\begin{enumerate}
    \item draw a sample of size \(n\) with replacement from the observed
          data;
    \item compute the statistic \(T^*\);
    \item repeat the process many times;
    \item use the empirical distribution of the bootstrap statistics to
          approximate the sampling distribution of \(T\).
\end{enumerate}

%%%%%%%%%%%%%%%%%%%%%%%%%%%%%%%%%%%%%%%%%%%%%%%%%%%%%%%%%%%%
\subsection{Bootstrap Standard Error}
\label{subsec:bootstrap-standard-error}
%%%%%%%%%%%%%%%%%%%%%%%%%%%%%%%%%%%%%%%%%%%%%%%%%%%%%%%%%%%%

Suppose bootstrap statistics are

\[
T_1^*,\ldots,T_B^*.
\]

Their sample standard deviation is

\[
\boxed{
\widehat{\operatorname{SE}}_{\text{boot}}
=
\sqrt{
\frac1{B-1}
\sum_{b=1}^{B}
(T_b^*-\overline{T^*})^2
}.
}
\]

This estimates the standard error of the original statistic.

%%%%%%%%%%%%%%%%%%%%%%%%%%%%%%%%%%%%%%%%%%%%%%%%%%%%%%%%%%%%
\subsection{Why Bootstrap Works Conceptually}
\label{subsec:why-bootstrap-works}
%%%%%%%%%%%%%%%%%%%%%%%%%%%%%%%%%%%%%%%%%%%%%%%%%%%%%%%%%%%%

The true sampling distribution depends on the unknown population
distribution \(F\).

The bootstrap replaces \(F\) by the empirical distribution

\[
F_n.
\]

It then simulates repeated sampling from \(F_n\).

Thus:

\[
\boxed{
F
\text{ unknown}
\longrightarrow
F_n
\text{ used as a plug-in approximation}.
}
\]

%%%%%%%%%%%%%%%%%%%%%%%%%%%%%%%%%%%%%%%%%%%%%%%%%%%%%%%%%%%%
\subsection{Bootstrap Limitations}
\label{subsec:bootstrap-limitations}
%%%%%%%%%%%%%%%%%%%%%%%%%%%%%%%%%%%%%%%%%%%%%%%%%%%%%%%%%%%%

Bootstrap procedures are powerful but not universally valid.

Problems may occur with:

\[
\boxed{
\text{very small samples},
}
\]

\[
\boxed{
\text{extreme-value statistics},
}
\]

\[
\boxed{
\text{strong dependence},
}
\]

\[
\boxed{
\text{irregular estimators},
}
\]

or

\[
\boxed{
\text{poorly representative samples}.
}
\]

Resampling cannot repair a fundamentally biased sampling design.

%%%%%%%%%%%%%%%%%%%%%%%%%%%%%%%%%%%%%%%%%%%%%%%%%%%%%%%%%%%%
\subsection{Monte Carlo Approximation of Sampling Distributions}
\label{subsec:monte-carlo-sampling-distributions}
%%%%%%%%%%%%%%%%%%%%%%%%%%%%%%%%%%%%%%%%%%%%%%%%%%%%%%%%%%%%

If the population model is known, a sampling distribution can be studied
by simulation.

For each simulation:

\begin{enumerate}
    \item generate a random sample from the assumed population model;
    \item compute the statistic;
    \item repeat many times.
\end{enumerate}

The resulting histogram approximates the sampling distribution.

This is different from bootstrap resampling because Monte Carlo
simulation may sample from a specified parametric model rather than the
empirical data.

%%%%%%%%%%%%%%%%%%%%%%%%%%%%%%%%%%%%%%%%%%%%%%%%%%%%%%%%%%%%
\subsection{Sampling Distribution and Repeated Sampling}
\label{subsec:repeated-sampling-interpretation}
%%%%%%%%%%%%%%%%%%%%%%%%%%%%%%%%%%%%%%%%%%%%%%%%%%%%%%%%%%%%

The phrase repeated sampling is conceptual.

It does not mean an investigator must physically repeat the study many
times.

Instead, the sampling distribution represents what would happen across
all hypothetical repetitions of the sampling procedure.

This probability model quantifies uncertainty in the one sample actually
observed.

%%%%%%%%%%%%%%%%%%%%%%%%%%%%%%%%%%%%%%%%%%%%%%%%%%%%%%%%%%%%
\subsection{Sampling Distribution and Confidence Intervals}
\label{subsec:sampling-confidence-connection}
%%%%%%%%%%%%%%%%%%%%%%%%%%%%%%%%%%%%%%%%%%%%%%%%%%%%%%%%%%%%

Suppose

\[
\widehat{\theta}
\]

has approximately normal sampling distribution

\[
\widehat{\theta}
\approx
N
(
\theta,
\operatorname{SE}^2
).
\]

Then

\[
\boxed{
\frac{
\widehat{\theta}-\theta
}{
\operatorname{SE}
}
\approx
N(0,1).
}
\]

This leads directly to confidence intervals of the form

\[
\boxed{
\widehat{\theta}
\pm
z^*
\operatorname{SE}.
}
\]

%%%%%%%%%%%%%%%%%%%%%%%%%%%%%%%%%%%%%%%%%%%%%%%%%%%%%%%%%%%%
\subsection{Sampling Distribution and Hypothesis Tests}
\label{subsec:sampling-testing-connection}
%%%%%%%%%%%%%%%%%%%%%%%%%%%%%%%%%%%%%%%%%%%%%%%%%%%%%%%%%%%%

A hypothesis test compares an observed statistic with its sampling
distribution under a null hypothesis.

If the observed statistic lies far into a tail of the null distribution,
the data may be considered inconsistent with that hypothesis.

Thus hypothesis testing depends fundamentally on:

\[
\boxed{
\text{the sampling distribution under }H_0.
}
\]

%%%%%%%%%%%%%%%%%%%%%%%%%%%%%%%%%%%%%%%%%%%%%%%%%%%%%%%%%%%%
\subsection{Exact Versus Approximate Sampling Distributions}
\label{subsec:exact-versus-approx-sampling}
%%%%%%%%%%%%%%%%%%%%%%%%%%%%%%%%%%%%%%%%%%%%%%%%%%%%%%%%%%%%

An exact sampling distribution is valid at the stated finite sample size.

Examples include:

\[
\boxed{
\overline{X}
\text{ from a normal population},
}
\]

\[
\boxed{
\frac{(n-1)S^2}{\sigma^2}
\sim
\chi^2_{n-1},
}
\]

and

\[
\boxed{
\frac{
\overline{X}-\mu
}{
S/\sqrt{n}
}
\sim
t_{n-1}
}
\]

under normal sampling.

An approximate sampling distribution is justified asymptotically, often
through the CLT.

%%%%%%%%%%%%%%%%%%%%%%%%%%%%%%%%%%%%%%%%%%%%%%%%%%%%%%%%%%%%
\subsection{Pivotal Quantities}
\label{subsec:pivotal-quantities}
%%%%%%%%%%%%%%%%%%%%%%%%%%%%%%%%%%%%%%%%%%%%%%%%%%%%%%%%%%%%

A pivotal quantity is a function of the data and unknown parameters whose
probability distribution does not depend on unknown parameters.

For normal data with known \(\sigma\),

\[
\boxed{
Z
=
\frac{
\overline{X}-\mu
}{
\sigma/\sqrt{n}
}
}
\]

has a standard normal distribution.

With unknown \(\sigma\),

\[
\boxed{
T
=
\frac{
\overline{X}-\mu
}{
S/\sqrt{n}
}
}
\]

has a \(t_{n-1}\) distribution.

Pivotal quantities are central in constructing confidence intervals.

%%%%%%%%%%%%%%%%%%%%%%%%%%%%%%%%%%%%%%%%%%%%%%%%%%%%%%%%%%%%
\subsection{Ancillary Statistics Preview}
\label{subsec:ancillary-statistics-preview}
%%%%%%%%%%%%%%%%%%%%%%%%%%%%%%%%%%%%%%%%%%%%%%%%%%%%%%%%%%%%

An ancillary statistic has a distribution that does not depend on the
parameter of interest.

Ancillary statistics can provide information about the sampling
configuration without directly estimating the parameter.

They become important in more advanced inference theory.

%%%%%%%%%%%%%%%%%%%%%%%%%%%%%%%%%%%%%%%%%%%%%%%%%%%%%%%%%%%%
\subsection{Sufficient Statistics Preview}
\label{subsec:sufficient-statistics-sampling-preview}
%%%%%%%%%%%%%%%%%%%%%%%%%%%%%%%%%%%%%%%%%%%%%%%%%%%%%%%%%%%%

A sufficient statistic contains all sample information relevant to a
parameter within a specified statistical model.

For example, for a normal model with known variance, the sample mean is
sufficient for the unknown mean.

Sufficiency is developed more fully in mathematical statistics.

%%%%%%%%%%%%%%%%%%%%%%%%%%%%%%%%%%%%%%%%%%%%%%%%%%%%%%%%%%%%
\subsection{Sampling Distribution Workflow}
\label{subsec:sampling-distribution-workflow}
%%%%%%%%%%%%%%%%%%%%%%%%%%%%%%%%%%%%%%%%%%%%%%%%%%%%%%%%%%%%

When analyzing a statistic:

\begin{enumerate}
    \item identify the population model or sampling design;
    \item identify the statistic;
    \item determine its expected value;
    \item determine its variance;
    \item compute its standard error;
    \item determine whether an exact distribution is available;
    \item otherwise determine whether a limit theorem applies;
    \item standardize the statistic when useful;
    \item check independence and sample-size assumptions;
    \item include finite-population corrections when appropriate;
    \item distinguish unknown parameters from their estimates;
    \item interpret the result as repeated-sampling behavior.
\end{enumerate}

%%%%%%%%%%%%%%%%%%%%%%%%%%%%%%%%%%%%%%%%%%%%%%%%%%%%%%%%%%%%
\subsection{Common Errors}
\label{subsec:sampling-distribution-common-errors}
%%%%%%%%%%%%%%%%%%%%%%%%%%%%%%%%%%%%%%%%%%%%%%%%%%%%%%%%%%%%

\subsubsection{Confusing the Population Distribution with the Sampling Distribution}

The population distribution describes individuals.

The sampling distribution describes a statistic.

They are not the same object.

%%%%%%%%%%%%%%%%%%%%%%%%%%%%%%%%%%%%%%%%%%%%%%%%%%%%%%%%%%%%
\subsubsection{Confusing a Statistic with Its Observed Value}

Before sampling,

\[
\overline{X}
\]

is random.

After sampling,

\[
\overline{x}
\]

is fixed.

%%%%%%%%%%%%%%%%%%%%%%%%%%%%%%%%%%%%%%%%%%%%%%%%%%%%%%%%%%%%
\subsubsection{Confusing Standard Deviation and Standard Error}

The standard deviation measures variation among observations.

The standard error measures variation among repeated statistics.

%%%%%%%%%%%%%%%%%%%%%%%%%%%%%%%%%%%%%%%%%%%%%%%%%%%%%%%%%%%%
\subsubsection{Using \(\sigma/\sqrt{n}\) for Dependent Data Without Adjustment}

The formula

\[
\frac{\sigma}{\sqrt{n}}
\]

assumes the relevant independent sampling structure.

Dependence can change the variance of the sample mean.

%%%%%%%%%%%%%%%%%%%%%%%%%%%%%%%%%%%%%%%%%%%%%%%%%%%%%%%%%%%%
\subsubsection{Ignoring Finite-Population Correction}

If a substantial fraction of a finite population is sampled without
replacement, the finite-population correction can materially reduce the
standard error.

%%%%%%%%%%%%%%%%%%%%%%%%%%%%%%%%%%%%%%%%%%%%%%%%%%%%%%%%%%%%
\subsubsection{Assuming Sample Means Are Always Exactly Normal}

Exact normality occurs when the underlying sample is normal.

For general populations, normality of the sample mean is usually an
approximation justified by the CLT.

%%%%%%%%%%%%%%%%%%%%%%%%%%%%%%%%%%%%%%%%%%%%%%%%%%%%%%%%%%%%
\subsubsection{Using the \(t\)-Distribution Without Its Assumptions}

The exact one-sample \(t\)-distribution requires normal sampling.

Large-sample procedures may be more robust, but the exact finite-sample
statement depends on the model.

%%%%%%%%%%%%%%%%%%%%%%%%%%%%%%%%%%%%%%%%%%%%%%%%%%%%%%%%%%%%
\subsubsection{Forgetting Degrees of Freedom}

For a normal sample variance,

\[
\boxed{
\frac{
(n-1)S^2
}{
\sigma^2
}
\sim
\chi^2_{n-1}.
}
\]

The degrees of freedom are \(n-1\), not \(n\).

%%%%%%%%%%%%%%%%%%%%%%%%%%%%%%%%%%%%%%%%%%%%%%%%%%%%%%%%%%%%
\subsubsection{Treating Paired Samples as Independent Samples}

Paired observations contain within-pair covariance.

Ignoring pairing can substantially change the standard error.

%%%%%%%%%%%%%%%%%%%%%%%%%%%%%%%%%%%%%%%%%%%%%%%%%%%%%%%%%%%%
\subsubsection{Using Estimated and Null Standard Errors Interchangeably}

Different inferential procedures may estimate standard errors under
different parameter assumptions.

For example, confidence intervals and null-hypothesis tests for
proportions can use different standard-error formulas.

%%%%%%%%%%%%%%%%%%%%%%%%%%%%%%%%%%%%%%%%%%%%%%%%%%%%%%%%%%%%
\subsubsection{Assuming Unbiased Means Accurate}

An estimator can be unbiased yet highly variable.

Bias alone does not determine estimator quality.

%%%%%%%%%%%%%%%%%%%%%%%%%%%%%%%%%%%%%%%%%%%%%%%%%%%%%%%%%%%%
\subsubsection{Assuming Consistent Means Unbiased}

A consistent estimator may have finite-sample bias that tends to zero as
sample size increases.

%%%%%%%%%%%%%%%%%%%%%%%%%%%%%%%%%%%%%%%%%%%%%%%%%%%%%%%%%%%%
\subsubsection{Assuming Bootstrap Resampling Fixes a Biased Sample}

Bootstrap methods reproduce the empirical sample distribution.

They do not automatically correct undercoverage, confounding, selection
bias, or nonresponse bias.

%%%%%%%%%%%%%%%%%%%%%%%%%%%%%%%%%%%%%%%%%%%%%%%%%%%%%%%%%%%%
\subsection{Applications}
\label{subsec:sampling-distribution-applications}
%%%%%%%%%%%%%%%%%%%%%%%%%%%%%%%%%%%%%%%%%%%%%%%%%%%%%%%%%%%%

\begin{applicationbox}

\textbf{Survey Research.}
Sampling distributions quantify uncertainty in estimates of population
means, proportions, and totals.

\medskip

\textbf{Scientific Experiments.}
Differences between treatment-group statistics are interpreted through
their sampling distributions.

\medskip

\textbf{Quality Control.}
Sampling distributions describe natural variation of process means,
proportions, and variances.

\medskip

\textbf{Clinical Studies.}
Treatment-effect estimates are standardized using their sampling
variability.

\medskip

\textbf{Education Research.}
Sample means, proportions, regression coefficients, and other statistics
vary across hypothetical samples of students.

\medskip

\textbf{Environmental Research.}
Sampling distributions quantify uncertainty in estimated concentrations,
means, trends, and group differences.

\medskip

\textbf{Data Science.}
Bootstrap and asymptotic sampling distributions provide uncertainty
estimates for complex statistics.

\medskip

\textbf{Statistical Inference.}
Confidence intervals and hypothesis tests are built directly from
sampling distributions.

\end{applicationbox}

%%%%%%%%%%%%%%%%%%%%%%%%%%%%%%%%%%%%%%%%%%%%%%%%%%%%%%%%%%%%
\subsection{Exercises}
\label{subsec:sampling-distribution-exercises}
%%%%%%%%%%%%%%%%%%%%%%%%%%%%%%%%%%%%%%%%%%%%%%%%%%%%%%%%%%%%

\begin{exercise}[1]
Define a sampling distribution.
\end{exercise}

\begin{exercise}[1]
Distinguish population, sample, and sampling distributions.
\end{exercise}

\begin{exercise}[1]
Distinguish an estimator from an estimate.
\end{exercise}

\begin{exercise}[1]
Define standard error.
\end{exercise}

\begin{exercise}[1]
State the mean and variance of the sample mean for an i.i.d. sample.
\end{exercise}

\begin{exercise}[1]
State the mean and variance of the sample proportion.
\end{exercise}

\begin{exercise}[1]
State the exact sampling distribution of the sample mean from a normal
population.
\end{exercise}

\begin{exercise}[1]
State the chi-square sampling distribution of the sample variance from a
normal population.
\end{exercise}

\begin{exercise}[1]
State the one-sample \(t\)-statistic and its degrees of freedom under
normal sampling.
\end{exercise}

\begin{exercise}[1]
Define bias.
\end{exercise}

\begin{exercise}[1]
Define mean squared error.
\end{exercise}

\begin{exercise}[1]
Define consistency.
\end{exercise}

\begin{exercise}[2]
Suppose

\[
\mu=50,
\qquad
\sigma=10,
\qquad
n=25.
\]

Find

\[
\mathbb{E}[\overline{X}]
\]

and

\[
\operatorname{SE}(\overline{X}).
\]
\end{exercise}

\begin{exercise}[2]
For the same sample, approximate

\[
P(\overline{X}>54)
\]

using a normal sampling distribution.
\end{exercise}

\begin{exercise}[2]
Suppose

\[
X_i\sim N(20,9).
\]

Find the exact distribution of the sample mean for

\[
n=16.
\]
\end{exercise}

\begin{exercise}[2]
Suppose

\[
p=0.30,
\qquad
n=100.
\]

Find the mean and standard error of

\[
\widehat{p}.
\]
\end{exercise}

\begin{exercise}[2]
Approximate

\[
P(\widehat{p}>0.35)
\]

for the previous problem.
\end{exercise}

\begin{exercise}[2]
Suppose

\[
N=500,
\qquad
n=100,
\qquad
\sigma=20.
\]

Compute the standard error of the SRS sample mean using the
finite-population correction.
\end{exercise}

\begin{exercise}[2]
Suppose

\[
n=10.
\]

What are the degrees of freedom of the sample variance chi-square
distribution?
\end{exercise}

\begin{exercise}[2]
Suppose

\[
n=20,
\qquad
\overline{x}=100,
\qquad
s=15,
\qquad
\mu_0=95.
\]

Compute the one-sample \(t\)-statistic.
\end{exercise}

\begin{exercise}[2]
Suppose two independent means have standard errors \(2\) and \(3\).
Find the standard error of their difference.
\end{exercise}

\begin{exercise}[2]
Suppose

\[
p_1=0.6,\quad n_1=100,
\qquad
p_2=0.4,\quad n_2=100.
\]

Find the standard error of

\[
\widehat{p}_1-\widehat{p}_2.
\]
\end{exercise}

\begin{exercise}[2]
For paired observations with

\[
\sigma_X=8,
\qquad
\sigma_Y=8,
\qquad
\rho=0.75,
\]

compute

\[
\operatorname{Var}(X-Y).
\]
\end{exercise}

\begin{exercise}[3]
Derive

\[
\mathbb{E}[\overline{X}]
=
\mu.
\]
\end{exercise}

\begin{exercise}[3]
Derive

\[
\operatorname{Var}(\overline{X})
=
\frac{\sigma^2}{n}
\]

under independence.
\end{exercise}

\begin{exercise}[3]
Derive the finite-population correction for an SRS sample mean.
\end{exercise}

\begin{exercise}[3]
Show that

\[
\widehat{p}
=
\frac Yn
\]

where

\[
Y\sim\operatorname{Bin}(n,p).
\]
\end{exercise}

\begin{exercise}[3]
Derive

\[
\operatorname{Var}(\widehat{p})
=
\frac{
p(1-p)
}{n}.
\]
\end{exercise}

\begin{exercise}[3]
Use the CLT to derive the asymptotic normal distribution of the sample
proportion.
\end{exercise}

\begin{exercise}[3]
Starting from the chi-square sampling distribution, prove that

\[
\mathbb{E}[S^2]=\sigma^2
\]

for normal samples.
\end{exercise}

\begin{exercise}[3]
Use the chi-square variance formula to derive

\[
\operatorname{Var}(S^2)
=
\frac{
2\sigma^4
}{
n-1
}
\]

for normal samples.
\end{exercise}

\begin{exercise}[3]
Explain why independence of \(\overline{X}\) and \(S^2\) is special to
normal sampling.
\end{exercise}

\begin{exercise}[3]
Derive the Student \(t\)-statistic from an independent standard normal
variable and chi-square variable.
\end{exercise}

\begin{exercise}[3]
Show that

\[
t_\nu
\xrightarrow{d}
N(0,1)
\]

as

\[
\nu\to\infty.
\]
\end{exercise}

\begin{exercise}[3]
Derive the standard error of the difference between two independent
sample means.
\end{exercise}

\begin{exercise}[3]
Derive the standard error of the difference between two independent
sample proportions.
\end{exercise}

\begin{exercise}[3]
Show that

\[
\operatorname{Var}(X-Y)
=
\operatorname{Var}(X)
+
\operatorname{Var}(Y)
-
2\operatorname{Cov}(X,Y).
\]
\end{exercise}

\begin{exercise}[3]
Explain mathematically why positive pairing correlation can reduce the
variance of a paired difference.
\end{exercise}

\begin{exercise}[3]
Derive the MSE decomposition

\[
\operatorname{MSE}
=
\operatorname{Var}
+
\operatorname{Bias}^2.
\]
\end{exercise}

\begin{exercise}[3]
Give an example of a biased but consistent estimator.
\end{exercise}

\begin{exercise}[3]
Give an example of two unbiased estimators with different variances.
\end{exercise}

\begin{exercise}[3]
Derive the CDF of the sample minimum.
\end{exercise}

\begin{exercise}[3]
Derive the CDF of the sample maximum.
\end{exercise}

\begin{exercise}[3]
Derive the density of the \(k\)-th order statistic.
\end{exercise}

\begin{exercise}[3]
Apply the delta method to

\[
g(\theta)=\theta^2.
\]
\end{exercise}

\begin{exercise}[3]
Explain conceptually how the bootstrap approximates a sampling
distribution.
\end{exercise}

\begin{exercise}[3]
Distinguish parametric Monte Carlo simulation from nonparametric
bootstrap resampling.
\end{exercise}

\begin{exercise}[4]
Prove the independence of sample mean and sample variance for a normal
sample.
\end{exercise}

\begin{exercise}[4]
Derive the exact chi-square distribution of the normal sample variance.
\end{exercise}

\begin{exercise}[4]
Derive the pooled two-sample \(t\)-distribution under equal normal
variances.
\end{exercise}

\begin{exercise}[4]
Study the Welch--Satterthwaite degrees-of-freedom approximation.
\end{exercise}

\begin{exercise}[4]
Study asymptotic normality of sample quantiles.
\end{exercise}

\begin{exercise}[4]
Investigate asymptotic distributions of M-estimators.
\end{exercise}

\begin{exercise}[4]
Study the multivariate delta method.
\end{exercise}

\begin{exercise}[4]
Investigate the bootstrap consistency theorem.
\end{exercise}

\begin{exercise}[4]
Study bootstrap failure for extreme-value estimators and irregular
statistics.
\end{exercise}

\begin{exercise}[4]
Investigate permutation distributions and compare them with sampling
distributions.
\end{exercise}

\begin{exercise}[4]
Study sufficient and ancillary statistics.
\end{exercise}

\begin{exercise}[4]
Investigate Basu's theorem.
\end{exercise}

\begin{exercise}[4]
Study asymptotic efficiency and Fisher information.
\end{exercise}

%%%%%%%%%%%%%%%%%%%%%%%%%%%%%%%%%%%%%%%%%%%%%%%%%%%%%%%%%%%%
\subsection{Conceptual Summary}
\label{subsec:sampling-distribution-summary}
%%%%%%%%%%%%%%%%%%%%%%%%%%%%%%%%%%%%%%%%%%%%%%%%%%%%%%%%%%%%

A statistic is a random variable before the sample is observed.

Its probability distribution under repeated sampling is its sampling
distribution.

For an i.i.d. sample,

\[
\boxed{
\mathbb{E}[\overline{X}]
=
\mu,
}
\]

\[
\boxed{
\operatorname{Var}(\overline{X})
=
\frac{\sigma^2}{n},
}
\]

and

\[
\boxed{
\operatorname{SE}(\overline{X})
=
\frac{\sigma}{\sqrt{n}}.
}
\]

For a sample proportion,

\[
\boxed{
\mathbb{E}[\widehat{p}]
=
p,
}
\]

and

\[
\boxed{
\operatorname{SE}(\widehat{p})
=
\sqrt{
\frac{
p(1-p)
}{n}
}.
}
\]

For normal samples,

\[
\boxed{
\overline{X}
\sim
N
\left(
\mu,
\frac{\sigma^2}{n}
\right),
}
\]

\[
\boxed{
\frac{
(n-1)S^2
}{
\sigma^2
}
\sim
\chi^2_{n-1},
}
\]

and

\[
\boxed{
\frac{
\overline{X}-\mu
}{
S/\sqrt{n}
}
\sim
t_{n-1}.
}
\]

For two independent means,

\[
\boxed{
\operatorname{SE}
(
\overline{X}-\overline{Y}
)
=
\sqrt{
\frac{
\sigma_X^2
}{
n_1
}
+
\frac{
\sigma_Y^2
}{
n_2
}
}.
}
\]

For two independent proportions,

\[
\boxed{
\operatorname{SE}
(
\widehat{p}_1-\widehat{p}_2
)
=
\sqrt{
\frac{
p_1(1-p_1)
}{
n_1
}
+
\frac{
p_2(1-p_2)
}{
n_2
}
}.
}
\]

Estimator quality can be described through

\[
\boxed{
\text{bias},
}
\]

\[
\boxed{
\text{variance},
}
\]

\[
\boxed{
\text{MSE},
}
\]

and

\[
\boxed{
\text{consistency}.
}
\]

Sampling distributions provide the probability framework that converts a
sample statistic into a measure of uncertainty.

%%%%%%%%%%%%%%%%%%%%%%%%%%%%%%%%%%%%%%%%%%%%%%%%%%%%%%%%%%%%
\subsection{Section Connections}
\label{subsec:sampling-distribution-connections}
%%%%%%%%%%%%%%%%%%%%%%%%%%%%%%%%%%%%%%%%%%%%%%%%%%%%%%%%%%%%

\chapterconnection{
Sampling Distributions formalize the sample-to-sample variability
introduced in Section~\ref{sec:data-collection-sampling}. The probability
theory of Chapter 7 determines the distributions of sample means,
proportions, variances, differences, order statistics, and other
statistics. Exact normal, chi-square, \(t\), and \(F\) distributions
provide the finite-sample foundation for classical statistical
inference, while the Central Limit Theorem provides asymptotic
approximations when exact formulas are unavailable. The next section,
Statistical Estimation, uses these sampling distributions to study point
estimators, bias, consistency, efficiency, likelihood, and methods for
estimating unknown population parameters. Confidence intervals and
hypothesis testing then use the same sampling-distribution framework to
quantify uncertainty and evidence.
}

%%%%%%%%%%%%%%%%%%%%%%%%%%%%%%%%%%%%%%%%%%%%%%%%%%%%%%%%%%%%
% SECTION SOURCE TRACKING
%%%%%%%%%%%%%%%%%%%%%%%%%%%%%%%%%%%%%%%%%%%%%%%%%%%%%%%%%%%%

%------------------------------------------------------------
% 8.3 Sampling Distributions
%------------------------------------------------------------
%
% Source basis:
%   - Standard mathematical statistics.
%   - Standard probability limit theory.
%   - Standard normal-theory sampling distributions.
%
% Used for:
%   - sampling-distribution definitions;
%   - population versus sample versus sampling distributions;
%   - estimators and estimates;
%   - sample-mean distributions;
%   - mean and variance of sample means;
%   - standard errors;
%   - finite-population corrections;
%   - exact normal sample means;
%   - central-limit approximations;
%   - sample sums;
%   - sample proportions;
%   - normal approximations to proportions;
%   - sample-variance distributions;
%   - chi-square distributions;
%   - independence of normal sample mean and variance;
%   - Student's t-distribution;
%   - F-distribution;
%   - ratios of sample variances;
%   - differences of independent means;
%   - pooled and Welch previews;
%   - differences of proportions;
%   - paired sampling;
%   - bias;
%   - estimator variance;
%   - mean squared error;
%   - consistency;
%   - asymptotic normality;
%   - delta-method preview;
%   - linear combinations;
%   - order-statistic distributions;
%   - minimum and maximum distributions;
%   - sample median preview;
%   - bootstrap sampling distributions;
%   - Monte Carlo approximation;
%   - pivotal quantities;
%   - sufficient and ancillary statistic previews;
%   - applications and exercises.
%
% Notes:
%   - Statistical Estimation is developed next in Section 8.4.
%   - Confidence Intervals are developed in Section 8.5.
%   - Hypothesis Testing is developed in Section 8.6.
%   - Normal-theory sampling distributions are reused throughout
%     regression, ANOVA, and mathematical statistics.
%
%%%%%%%%%%%%%%%%%%%%%%%%%%%%%%%%%%%%%%%%%%%%%%%%%%%%%%%%%%%%